% 《大学物理简明教程》（第5版）第4章选择题 4.1.1—4.1.12。
% 书内页/PDF页100—102；依据扫描题干和图形，非官方改述+原创推导。
% 自检不等于独立学术审定。
\subsection{选择题：简谐运动、波及干涉}

\sourceproblem{4.1.1}{100}{100}
\problemtext 一质点做简谐运动，$t=0$时$x=-A/2$且向$x$正方向运动；从旋转矢量法判断初相位。
\workings
以$x=A\cos(\omega t+\varphi)$表示，初始条件给
$\cos\varphi=-1/2$，而
$v(0)=-A\omega\sin\varphi>0$，所以$\sin\varphi<0$。
\[
\boxed{\varphi=\frac{4\pi}{3}\pmod{2\pi}}
\]
旋转矢量位于第三象限（与图示对应的选项为B）。
\crosscheck 若误取$2\pi/3$，速度会指向负$x$，不满足题设方向。

\sourceproblem{4.1.2}{100}{100}
\problemtext 两同周期简谐振动的位移—时间曲线如原图4.1.2，曲线(1)初始位于负最大位移，而曲线(2)先于(1)达到正最大位移。判断(1)相对于(2)的相位。
\workings
观察同一正峰位置，曲线(1)向右滞后四分之一周期$\Delta t=T/4$，故
\[
\boxed{\varphi_1-\varphi_2=-\omega\Delta t=-\frac{\pi}{2}}
\]
（模$2\pi$），即(1)落后(2) $\pi/2$，原题选A。
\crosscheck 对同频曲线使用峰间\emph{水平时间差}，不是直接比较不同峰的位移高低。

\sourceproblem{4.1.3}{100--101}{100--101}
\problemtext 做简谐运动的质点，周期为$T$；在$x=A/2$处向$x$正方向运动，问此后第一次到达$x=A$需要多长时间。
\workings
由$\cos\varphi=1/2$且$v=-A\omega\sin\varphi>0$，
取$\varphi=-\pi/3$。到达正极端的相位为$0$，变化$\Delta\theta=\pi/3$，故
\[
\boxed{\Delta t=\frac{\pi/3}{2\pi/T}=\frac T6}
\]
（原题选C）。
\crosscheck 非线性匀速位移法$A/2v$不适用，必须以相位差计算。

\sourceproblem{4.1.4}{101}{101}
\problemtext 连接同一轻弹簧的物体置于竖直方向或光滑斜面上，在弹簧始终服从胡克定律且物体受适当约束时，能否做简谐振动？
\workings
沿允许运动的坐标$s$，重力分力为常数$F_g$，弹力为$-ks$，
\[
m\ddot s=F_g-ks.
\]
设平衡位置$s_e=F_g/k$，令$\xi=s-s_e$，便有
$m\ddot\xi+k\xi=0$，即$\xi$为简谐运动。竖直和光滑斜面两情形均成立（选C）。
\crosscheck 这以弹簧保持连接、无松弛、无碰撞和线性胡克定律为前提；重力只是移动平衡位置，不改变$\omega=\sqrt{k/m}$。

\sourceproblem{4.1.5}{101}{101}
\problemtext $x=A\cos(\omega t+\varphi)$，在$t=T/2$时瞬时速度是多少？
\workings
$\omega T=2\pi$，
\[
v(T/2)=-A\omega\sin(\pi+\varphi)
=\boxed{A\omega\sin\varphi}.
\]
原题选B。
\crosscheck 半周期后速度与初速反号；初速$-A\omega\sin\varphi$。

\sourceproblem{4.1.6}{101}{101}
\problemtext 质点以频率$\nu$作简谐振动，问动能随时间的变化频率。
\workings
由$v=-A\omega\sin(\omega t+\varphi)$，
\[
K=\frac12mA^2\omega^2\sin^2(\omega t+\varphi)
=\frac14mA^2\omega^2[1-\cos(2\omega t+2\varphi)].
\]
因此能量变化角频率是$2\omega$，频率
$\boxed{\nu_K=2\nu}$，原题选C。
\crosscheck 质点经过一个周期有两次速度最大、两次速度为零。

\sourceproblem{4.1.7}{101}{101}
\problemtext 在简谐振动中，动能与势能相等的位置是什么？
\workings
总能量$E=\frac12kA^2$，势能$U=\frac12kx^2$，
$K=E-U$。令$K=U$，得
\[
kx^2=\frac12kA^2,\qquad
\boxed{x=\pm\frac{A}{\sqrt2}}.
\]
原题选D。
\crosscheck 两位置关于平衡点对称，能量均为总能量一半。

\sourceproblem{4.1.8}{101}{101}
\problemtext 弹簧振子质量、弹簧劲度系数不变，振幅变为原来的两倍。比较周期、能量、最大速率与最大加速度。
\workings
\[
T=2\pi\sqrt{\frac mk},\quad
E=\frac12kA^2,\quad
v_{\max}=\omega A,\quad
a_{\max}=\omega^2 A.
\]
令$A\mapsto2A$后，$T$不变，$E$变为$4E$，
$v_{\max}$及$a_{\max}$均变为2倍。原题选C。
\crosscheck 线性弹簧振子中周期与振幅无关，不适用于任意大幅非线性振动。

\sourceproblem{4.1.9}{101}{101}
\problemtext 波沿$x$正方向以速率$u$传播；距$O$点$l$的$P$点满足$y_P=A\cos\omega t$，求$O$点和整条波的波函数。
\workings
$P$在$x=l$，正向行波一般写为$y(x,t)=A\cos[\omega t-k(x-l)]$，$k=\omega/u$。
因此
\[
\boxed{y(x,t)=A\cos\omega\left(t+\frac lu-\frac xu\right)},\qquad
\boxed{y_O(t)=A\cos\omega\left(t+\frac lu\right)}.
\]
原题对应选C。
\crosscheck 代入$x=l$恢复$P$点已知振动；固定相位条件$x-ut=\mathrm{const}$为正向传播。

\sourceproblem{4.1.10}{101}{101}
\problemtext 一机械波$y=0.05\cos(6\pi t+0.06\pi x)$（SI）。判断波长、波速、周期与传播方向。
\workings
由$\omega=6\pi\,\mathrm{rad/s}$、$k=0.06\pi\,\mathrm{rad/m}$可得
\[
f=3\,\mathrm{Hz},\quad
\boxed{T=\frac13\,\mathrm s},\quad
\lambda=\frac{2\pi}{k}=\frac{100}{3}\,\mathrm m,\quad
u=\frac{\omega}{k}=100\,\mathrm{m/s}.
\]
相位$\omega t+kx$为常量时$x$随$t$减小，故沿$x$负方向传播。
四个选项中仅周期为$1/3\,\mathrm s$正确（C）。
\crosscheck 不能将方程中$0.06\pi$直接当成$2\pi/\lambda$之外的“波长”。

\sourceproblem{4.1.11}{101--102}{101--102}
\problemtext 两相干波源$S_1,S_2$相位分别为$\varphi_1,\varphi_2$，到$P$距离$r_1,r_2$，波长均为$\lambda$。写出在$P$干涉加强的相位条件。
\workings
两波到达$P$的相位分别为
$\theta_1=\omega t+\varphi_1-2\pi r_1/\lambda$与
$\theta_2=\omega t+\varphi_2-2\pi r_2/\lambda$。
相位差
\[
\Delta\theta=\theta_2-\theta_1
=\varphi_2-\varphi_1+\frac{2\pi(r_1-r_2)}{\lambda}.
\]
因此干涉加强的条件是
\[
\boxed{\varphi_2-\varphi_1+\frac{2\pi(r_1-r_2)}{\lambda}=2n\pi},
\quad n\in\mathbb Z.
\]
原题选D。
\crosscheck $r_1=r_2$且$\varphi_1=\varphi_2$时必加强；若两源反相需半整数波长差。

\sourceproblem{4.1.12}{102}{102}
\problemtext 弹性介质中传播的平面简谐行波，一质元从平衡位置移动到最大位移位置的过程中，它的机械能如何变化？
\workings
对理想线性无耗散行波，局部动能密度和弹性势能密度同相变化。以
$y=A\cos(\omega t-kx)$为例，
\[
\mathcal E_k=\frac12\rho(\partial_t y)^2
=\frac12\rho A^2\omega^2\sin^2(\omega t-kx),
\]
而对一维弦的弹性势能密度
$\mathcal E_p=\frac12T(\partial_x y)^2
=\frac12\rho A^2\omega^2\sin^2(\omega t-kx)$
（用$T/\rho=u^2$适当单位化）。
质元由平衡位移达到最大位移时，两个局部能量密度均由最大变为零，
能量通过介质传给相邻部分。原题选D。
\crosscheck 传播波的一小段介质并非孤立弹簧振子，局部总能量不必守恒；静止正弦驻波的能量分布则不同。
